% Paper 2 (ReSIDS v2) -- limitations (discussion section). Impersonal style.
% Evidence: docs/RESULTADOS_ciciot2023.md (label/protocol check), docs/RESULTADOS_electra.md
% (path labels, process transient), docs/paper_v2_dataset_audit.tex (Morris-4, ERENO).
% Decided 2026-10-07: context-aware detection (per-source state, benign model of
% communication relations, Tier-2 specification rules) stays OUT of Paper 2; cited as
% future work only.

\paragraph{Limitations.}
The specialists judge one sample at a time, from its content. Attacks defined by their
context rather than their content therefore escape them: the spoofing and
man-in-the-middle attacks of CICIoT2023, whose ARP packets never reach the flow records;
the replayed packets and the read requests relayed by the attacker's host in Electra,
which carry the content of normal packets; and the masquerade attacks of ERENO, which
mimic the legitimate publisher. Detecting them requires state kept per source, such as
rates, sequence numbers and repeated content, or a model of the communication relations
learned from benign traffic alone, in which a new path is itself the anomaly. Neither is
part of this work. The specialists are also supervised, and they reproduce the errors of
their labels: in CICIoT2023, labels that mark the capture session taught them to flag
benign traffic as attack, and nothing inside ReSIDS revealed it. The external audit did,
prompted by the false-positive rate that the evaluation at realistic prevalence exposed.
Process transients are a third source of error: in Electra, seven minutes of legitimate
register values outside their usual range were taken for modified responses. Finally, the IoT
domain rests on one dataset and three attack categories. Of the eight public IoT and
industrial IoT datasets audited, only the UNSW IoT traces support a sound evaluation, and
there, at one-minute resolution, TCP and UDP floods against busy devices fall within normal
variation; reconnaissance, web and brute-force attacks on IoT devices remain without a
dataset that would allow their evaluation. ARP spoofing, by contrast, is detected there:
counters aggregated per device and minute already carry the context that a single packet
lacks.
